%==============================================================================
% Predicting Quarterly Airbnb Booking Counts in New York City
% Journal manuscript -- IEEEtran source (double-blind)
%
% Compile with:  pdflatex main  ->  bibtex main  ->  pdflatex main  ->  pdflatex main
% (or upload the whole paper/ folder to Overleaf and press Recompile)
%
% This revision addresses the peer-review feedback by adding:
%   * formal equations for the K-fold target encoder, MLP, randomized search
%     with early stopping, and the SLSQP-constrained blend (Sec. IV);
%   * a mathematical critique of linear models on zero-inflated targets (Sec. II);
%   * a statistical-significance test for the blend gain and a residual
%     analysis on the [0,92] boundaries (Sec. V);
%   * DOIs for the journal references (references.bib).
%==============================================================================
\documentclass[journal]{IEEEtran}

\usepackage{amsmath,amssymb,amsfonts}
\usepackage{graphicx}
\usepackage{textcomp}
\usepackage{booktabs}
\usepackage{url}
\usepackage[T1]{fontenc}
\usepackage[utf8]{inputenc}
\usepackage{xcolor}
\graphicspath{{figures/}}

\hyphenation{op-tical net-works semi-conduc-tor}

\begin{document}

\title{Predicting Quarterly Airbnb Booking Counts in New York City with
Tree-Based Ensembles and a Weighted Model Blend}

\author{Anonymous Author(s)%
\thanks{Author names and affiliations omitted for double-blind review.}}

\markboth{Manuscript submitted for review}%
{Predicting Quarterly Airbnb Booking Counts in New York City}

\maketitle

\begin{abstract}
Forecasting how heavily a short-term rental will be booked is a valuable but
difficult problem for both hosts and the platforms that connect them with
guests. This study addresses the prediction of the total number of nights booked
in the third quarter of 2016 for approximately 24{,}000 short-term rental
properties in New York City, evaluated by Mean Squared Error (MSE) on a held-out
test set. The target is a censored, zero-inflated count in $[0,92]$, with a
$53.5\%$ mass at zero, which makes it a challenging demand-forecasting task.
Starting from five raw tables (one static property table and two quarters of
daily listing histories totalling more than eight million rows), a flat matrix
of 150 features per property is engineered, covering price dynamics, availability
and booking behaviour, status churn, recency, a quarter-over-quarter demand
proxy, text keywords from listing titles, and geographic clusters. Five
regressors---Linear Regression, Random Forest, Gradient Boosting, XGBoost, and a
PyTorch Multi-Layer Perceptron (MLP) trained with Adam---are compared inside
leakage-safe scikit-learn pipelines, and the boosted trees are tuned with a
randomized search that uses fold-level early stopping. The best single model, a
tuned XGBoost regressor, reaches a 5-fold cross-validation MSE of 182.2, and a
convex (SLSQP) weighted blend lowers the out-of-fold MSE further to 180.0. The
target-encoding, search, MLP, and blending objectives are defined formally, and
the blend gain is verified with a paired significance test
($t=5.33$, $p=1.0\times10^{-7}$) and a residual analysis on the bounded $[0,92]$
target. The proposed blend also outperforms strong gradient-boosting baselines
(LightGBM, 184.1; CatBoost, 184.3) on the same split, and an ablation study
quantifies the contribution of each pipeline component.
\end{abstract}

\begin{IEEEkeywords}
Airbnb, demand forecasting, regression, ensemble learning, XGBoost, random
forest, multi-layer perceptron, model blending, target encoding, constrained
optimization, mean squared error, feature engineering.
\end{IEEEkeywords}

%==============================================================================
\section{Introduction}
\IEEEPARstart{O}{ver} the last ten years, peer-to-peer (P2P) accommodation
platforms have reshaped the way people travel and the way property owners earn
money from their homes. Airbnb is the most visible example of this shift, with
millions of active listings spread across more than a hundred thousand cities
worldwide \cite{adamiak2018}. For a platform of that size, two questions tend
to dominate operations: what is a fair price for a given listing, and how much
demand will that listing receive over a future period. This study addresses the
second question and carries it through to a complete, empirically validated
solution.

Predicting demand a quarter in advance is valuable for several practical
reasons. A host who anticipates a busy three-month period can schedule cleaning
and maintenance, raise prices for high-demand weeks, or hire help. The platform
can plan customer-support capacity, target marketing toward listings that appear
likely to be under-booked, and provide better pricing recommendations.
Predicting demand remains challenging, however, because it arises from the
complex interplay of many factors---property type and capacity, neighbourhood,
seasonality, host responsiveness, review history, and how aggressively the price
moves during the period. Classical linear models have been applied to such
accommodation data \cite{gunter2018,wang2017}, but a consistent finding in the
recent literature is that machine-learning models, and tree-based ensembles in
particular, perform better on the mixed numerical and categorical tables that
describe short-term rentals \cite{tang2024,camatti2024,kalehbasti2019}.

The task studied here is to predict, for every property in a held-out New York
City test set, the total number of nights booked during the third quarter (July,
August, and September) of 2016. Because a quarter contains 92 days, the target is
an integer in $[0,92]$, and predictions are evaluated by Mean Squared Error
(MSE). Two structural properties make this harder than the price-prediction
tasks that dominate the literature: the target is \emph{censored} at both ends
and is \emph{zero-inflated}, with $53.5\%$ of properties booked for zero nights
in the quarter.

\emph{Research gap.} While the current literature applies tree-based models
extensively to short-term-rental \emph{price} prediction, systematic frameworks
that jointly address extreme zero-inflation (here a $53.5\%$ zero mass) and a
hard bounded support $[0,92]$ in a \emph{demand}-forecasting setting, and that
combine heterogeneous learners through an optimized convex ensemble, remain
under-explored. This work targets that gap: it pairs leakage-safe temporal
feature engineering with a constrained (SLSQP) convex blend, formalizes each
component mathematically, and validates the resulting gain statistically and
against modern gradient-boosting baselines.

The remainder of the paper is organized as follows. Section~\ref{sec:lit}
reviews related work on rental-demand prediction, tree-based ensembles,
hyper-parameter tuning, and preprocessing, and gives a mathematical argument for
why linear models struggle with a zero-inflated bounded target.
Section~\ref{sec:settings} states the problem and describes the dataset.
Section~\ref{sec:method} presents the full methodology, defining every component
---the target encoder, the search procedure, the MLP, and the blend---with
explicit equations. Section~\ref{sec:results} reports the results, the
comparison against SOTA baselines, a significance test, a residual analysis, and
a failure-case study. Section~\ref{sec:changes} presents an ablation study.
Section~\ref{sec:concl} discusses limitations and concludes.

%==============================================================================
\section{Literature Review}
\label{sec:lit}
The related work is grouped into four areas: (A) machine learning for Airbnb
prediction tasks, (B) decision trees and tree-based ensembles, (C)
hyper-parameter tuning and model validation, and (D) data preprocessing and
feature engineering.

\subsection{Machine Learning for Airbnb Prediction Tasks}
Several studies apply machine learning directly to Airbnb data. Tang, Cheng, and
Zhang \cite{tang2024} worked with listings from Sydney and compared ten
algorithms including Random Forest, Gradient Boosting, and CatBoost for price
prediction. CatBoost gave them the lowest root mean squared error, and the most
influential features were the maximum number of guests, the number of bedrooms,
and whether the room was private or shared. Camatti et al.\ \cite{camatti2024}
ran a comparative analysis of artificial-intelligence and traditional linear
methods on Dutch Airbnb listings, and concluded that tree-based and
neural-network methods overfit less than linear regression, though linear models
are still handy for reading off which predictors matter. Rezazadeh Kalehbasti,
Nikolenko, and Rezaei \cite{kalehbasti2019} predicted New York City Airbnb
prices using property features, host characteristics, and the sentiment of
customer reviews; their best model was a regularized linear model built on
review-text features, but tree models such as Random Forest and Support Vector
Regression were close behind.

Other work looks at demand rather than price, which is closer to the present task.
Siker \cite{siker2018} tackled the Kaggle ``Airbnb New User Bookings'' problem
with XGBoost, and even though the target there (destination country) is
different from ours, the workflow---heavy feature engineering followed by
gradient boosting---matches how most Airbnb-style Kaggle problems are solved in
practice. Chapman, Mohammad, and Villegas \cite{chapman2023} studied dynamic
short-term rental markets and built a pipeline that merged listing, calendar,
and review data; they again found gradient boosting beating linear regression.
Islam et al.\ \cite{islam2022} combined a Latent Dirichlet Allocation topic
model over property descriptions with an XGBoost regressor and improved price
prediction on the San Jose dataset. Two lessons from this group informed the
present design: first, the most useful predictors are usually a mix of structural
features, location, and dynamic calendar/review information; and second,
gradient-boosted trees are a strong default because they cope well with mixed
feature types and missing values.

A recurring theme across these studies is that the engineering choices around the
model often matter more than the model itself. Tang et al.\ \cite{tang2024} and
Islam et al.\ \cite{islam2022} both report that their largest gains came from the
features they constructed---guest capacity and room type in the first case, topic
structure of the listing description in the second---rather than from swapping one
boosting library for another. Kalehbasti et al.\ \cite{kalehbasti2019} similarly
found that review-sentiment features moved their error more than the choice
between a regularized linear model and a tree ensemble. This is consistent with
the broader observation that on tabular problems the ceiling is usually set by how
much signal the practitioner can encode, and it is precisely why most effort
went into the 150-feature matrix, with the model comparison treated as a
secondary, if still important, step. It also explains why a competition with a
fixed, shared dataset is decided largely by feature engineering and validation
discipline rather than by access to a particular algorithm---an observation that
shaped the allocation of effort.

\subsection{Why Linear Models Underfit a Zero-Inflated Bounded Target}
\label{sec:linearcritique}
Beyond the empirical observation that linear regression trails the trees, it is
worth stating \emph{why} in mathematical terms, because the shape of the target
is the crux of the whole problem. Ordinary least squares fits
$\hat y(\mathbf{x}) = \boldsymbol{\beta}^{\top}\mathbf{x}$ by minimizing the
population risk $\mathbb{E}\big[(y-\boldsymbol{\beta}^{\top}\mathbf{x})^{2}\big]$,
whose minimizer is the conditional mean $\mathbb{E}[y\mid\mathbf{x}]$
\emph{restricted to the class of affine functions}. The target, however, is a
censored count confined to $[0,92]$ with a large probability mass at $0$
(roughly 54\% of the development rows are exactly zero; see
Fig.~\ref{fig:target}). Its true conditional mean is therefore an inherently
nonlinear, saturating function of the covariates: it must flatten to $0$ for
low-demand listings and saturate toward $92$ for high-demand ones. An affine
predictor cannot reproduce those two plateaus---it extrapolates linearly past
both edges, so it produces negative fitted values for the dead listings (which
are then clipped, wasting capacity) and undershoots the busiest ones. Formally,
if $y=\max\!\big(0,\min(92,\ \boldsymbol{\beta}^{\top}\mathbf{x}+\eta)\big)$ with
a zero-mean noise term $\eta\sim\mathcal{N}(0,\sigma^{2})$,
the censoring makes $\mathbb{E}[y\mid\mathbf{x}]$ a \emph{nonlinear} function of
$\mathbf{x}$ even when the latent variable is linear, which is exactly the
limited-dependent-variable setting that motivated Tobit-type estimators
\cite{tobin1958} and, more directly for counts, two-part (hurdle) models.
Earlier price-prediction studies that reported linear baselines
\cite{gunter2018,wang2017} did not face this point mass because price is
continuous and unbounded; demand counts are not, and that is the structural
reason their linear formulations do not transfer to demand counts. Tree ensembles
sidestep the issue entirely: a regression tree approximates the saturating
conditional mean with axis-aligned piecewise-constant regions, so the $0$ and
$92$ plateaus are simply leaves, and boosting refines the transition region
between them.

\subsection{Modelling Zero-Inflated and Count Targets}
\label{sec:zeroinflated}
The target is a bounded count with a large spike at zero, and the statistics
literature offers a dedicated family of models for exactly this shape. A plain
Poisson regression assumes the conditional mean equals the conditional variance,
which is violated here because the zero spike and the upper-bound pile-up make
the data strongly over-dispersed. Negative-binomial regression relaxes that
equal-mean-variance assumption but still does not account for the excess of
structural zeros. Two model classes target the zeros directly: zero-inflated
models (e.g.\ zero-inflated Poisson, ZIP) mix a point mass at zero with a count
distribution, and \emph{hurdle} (two-part) models first decide whether the count
is positive and then model the positive counts separately. Censored-regression
estimators such as Tobit \cite{tobin1958} instead treat the bound itself as the
generative mechanism. These parametric estimators are not adopted in the final
pipeline---they are outside the scope of this study, and gradient-boosted trees
already approximate the saturating conditional mean non-parametrically---but they
frame the problem precisely and directly motivate the two-stage hurdle extension
proposed in Section~\ref{sec:concl}. The residual analysis in
Section~\ref{sec:residuals} shows empirically why a single least-squares
regressor leaves room for such a two-part treatment: its errors are systematically
signed at both the zero spike and the upper bound.

\subsection{Decision Trees and Tree-Based Ensembles}
Tree-based methods are central to the proposed solution. The decision tree, and Quinlan's well-known ID3 variant
\cite{quinlan1986}, showed that recursively partitioning the feature space
yields interpretable, non-parametric classifiers and regressors. A single tree
overfits easily, so two families of ensembles were developed to fix that.
Bagging \cite{breiman1996} trains many trees on bootstrap samples and averages
their predictions, and Random Forests \cite{breiman2001} extend bagging by also
sampling a random subset of features at each split, which decorrelates the trees
and cuts variance. Boosting, and Gradient Boosting in particular
\cite{friedman2001}, instead builds trees sequentially so that each new tree
fits the residual errors of the current ensemble. XGBoost \cite{chen2016xgboost}
is a heavily optimized open-source implementation of gradient boosting with
sparsity-aware split finding, $L_1/L_2$ regularization, and early stopping---all
features that turned out to matter for the missing-value-heavy Airbnb tables.
The textbook by Hastie, Tibshirani, and Friedman \cite{hastie2009esl} gives the
unified statistical picture behind all of these, and scikit-learn
\cite{pedregosa2011sklearn} is used as the common Python interface throughout.

\subsection{Hyper-Parameter Tuning and Model Validation}
Tree ensembles have several knobs---number of trees, learning rate, maximum
depth, subsampling rates, and regularization terms. Bergstra and Bengio
\cite{bergstra2012random} showed that random search beats grid search for
hyper-parameter optimization when only a few of those knobs really drive the
score, which is usually the case. This advice is followed; however, as
Section~\ref{sec:changes} explains, a custom randomized
search loop instead of using scikit-learn's \texttt{RandomizedSearchCV},
specifically to combine the search with fold-level early stopping.
Cross-validation is what makes the model comparison trustworthy; it
lowers the variance of the performance estimate and prevents tuning to a
single lucky split. For comparing two trained models on the same data,
Dietterich \cite{dietterich1998} recommends paired resampling tests, which are
adopted in Section~\ref{sec:significance}.

\subsection{Data Preprocessing, Feature Engineering, and Ensembling}
Beyond the model, careful preprocessing has a large effect on the final number.
The usual steps---imputing missing values, encoding categoricals, scaling
numerics, and, above all, preventing leakage by wrapping everything in
\texttt{Pipeline}s and \texttt{ColumnTransformer}s
\cite{pedregosa2011sklearn}---constitute standard leakage-safe practice. For
high-cardinality categorical columns, $K$-fold target encoding is used, a
leakage-aware version of mean encoding, defined formally in
Section~\ref{sec:preproc}. Finally, the proposed solution does not rely on a
single estimator: it blends several models with non-negative weights that sum to
one, an idea that goes back to Wolpert's stacked generalization
\cite{wolpert1992stacked}. The MLP component uses PyTorch \cite{paszke2019pytorch}
with the Adam optimizer \cite{kingma2015adam}.

%==============================================================================
\section{Project Settings}
\label{sec:settings}
\subsection{Problem Statement and Objectives}
The problem is to predict, for each Airbnb property in the test set, the total
number of nights booked in the third quarter of 2016 (July, August, September)
in New York City. Formally, given a feature vector $\mathbf{x}_i$ describing
property $i$, the goal is a function $f$ such that
$\hat y_i = f(\mathbf{x}_i)$, where the true target $y_i$ is a non-negative
integer in the range $[0,92]$ (92 days in Q3). The competition metric is the
Mean Squared Error,
\begin{equation}
\text{MSE} = \frac{1}{n}\sum_{i=1}^{n}\big(y_i-\hat y_i\big)^{2},
\label{eq:mse}
\end{equation}
so a lower MSE is better.

The concrete objectives are: (i) to build a fully reproducible, leakage-safe
pipeline; (ii) to clearly beat a constant-mean baseline and a linear baseline;
(iii) to interpret the most important features so the predictions are not a
black box; and (iv) to validate the proposed ensemble against strong
gradient-boosting baselines on a held-out test set. As reported in
Section~\ref{sec:results}, all four are met.

\subsection{Dataset}
The data follow the structure popularised by large-scale public short-term-rental
scrapes of New York City (of the kind distributed by the Inside Airbnb project
and used widely in the literature \cite{kalehbasti2019}). They comprise five
tables. A static property table holds information for 48{,}690 properties (33
columns: property and listing type, neighbourhood, zip code, capacity, published
nightly/weekly/monthly rates, fees, review counts and ratings, host flags,
latitude/longitude, and creation date). Two daily listing tables give the daily
status and price of each listing for the first and second quarters of 2016
---about 4.19 and 4.40 million rows respectively---with a \texttt{Status} field
taking the values A (available), B (booked), and R (reserved/blocked). A training
label table provides Q3 booking counts for 24{,}372 properties, and a test table
lists the 24{,}318 properties to be predicted.

This dataset is well suited to the demand-forecasting problem for three reasons.
First, it is real New York City short-term-rental data, one of the busiest
markets in the world \cite{kalehbasti2019}. Second, it poses a well-defined
regression task with a clearly specified metric and a held-out test split.
Third, and most distinctively, it requires combining a static property table
with two long daily histories; as shown later, the daily files carry most of the
predictive signal.

%==============================================================================
\section{Methodology}
\label{sec:method}
This section describes the final pipeline end to end. The order matches the
actual processing stages: exploratory analysis, feature engineering,
the four baseline model families, the boosted-tree retuning, and finally the
blend. In response to reviewer feedback, every custom component is now stated as
an explicit optimization problem rather than only described in prose.

\subsection{Exploratory Data Analysis}
Prior to feature extraction, an exploratory data analysis (EDA) was conducted. Three
observations shaped later decisions. First, the target is heavily skewed
and zero-inflated: a large group of properties is booked for very few or zero
nights in Q3, while a smaller group is booked almost the whole quarter
(Fig.~\ref{fig:target}). Second, the static table is full of
holes---\texttt{ExtraPeopleFee} is missing for 59\% of rows, \texttt{Security\-Deposit}
for 53\%, \texttt{ResponseRate} for 40\%, \texttt{CleaningFee} for 29\%, and
\texttt{OverallRating} for 25\%---so missingness itself is probably informative,
not just noise. Third, when simple per-property aggregates were computed from the
daily files and correlated them with the target, the strongest single predictors
were all behavioural: the reserved-day count and reserved rate in Q2, the number
of unique reservations, and the number of reviews per month (all with
correlations around 0.70--0.78). That indicated early that the daily files, not
the static table, were where the signal lived.

\begin{figure}[!t]
\centering
\includegraphics[width=\columnwidth]{fig1_target_dist.png}
\caption{Distribution of the target (Q3 2016 booked nights). The target is
zero-inflated on the left and piles up again near the 92-day maximum, which
later justified clipping predictions to $[0,92]$.}
\label{fig:target}
\end{figure}

Table~\ref{tab:targetstats} quantifies the shape of the target on the
development split. The mean (18.55) sits far above the median (0), the
distribution is strongly right-skewed (skewness $1.33$), and---most
importantly---53.5\% of properties were booked for zero nights in Q3, while only
about 1\% reach the upper bound. This single statistic is what makes the task a
\emph{zero-inflated, bounded} regression problem rather than an ordinary one,
and it is the structural reason a least-squares estimator is biased at the
extremes (Sections~\ref{sec:linearcritique} and~\ref{sec:residuals}).
Table~\ref{tab:missing} lists the static columns with the highest missing rates;
because these are far from missing-at-random (a missing \texttt{ResponseRate}
typically marks an inactive host), missingness is encoded explicitly with
indicator columns rather than discarding it.

\begin{table}[!t]
\caption{Development-set target statistics ($n=19{,}497$).}
\label{tab:targetstats}
\centering
\begin{tabular}{lc}
\toprule
Statistic & Value \\
\midrule
Mean & 18.55 \\
Median & 0 \\
Standard deviation & 27.80 \\
Skewness & 1.33 \\
Rows with $y=0$ & 53.5\% \\
Rows with $y\ge 90$ & 1.1\% \\
Rows with $y=92$ & 0.5\% \\
\bottomrule
\end{tabular}
\end{table}

\begin{table}[!t]
\caption{Missing-value rates of the most affected static columns.}
\label{tab:missing}
\centering
\begin{tabular}{lc}
\toprule
Column & Missing rate \\
\midrule
\texttt{ExtraPeopleFee} & 59\% \\
\texttt{SecurityDeposit} & 53\% \\
\texttt{ResponseRate} & 40\% \\
\texttt{CleaningFee} & 29\% \\
\texttt{OverallRating} & 25\% \\
\bottomrule
\end{tabular}
\end{table}

\subsection{Feature Engineering}
The core of the project was turning the three raw tables into a single flat
matrix with one row per \texttt{PropertyID}. The result is 150 features (144
numeric and 6 categorical). They fall into the following groups.
\begin{itemize}
\item \textbf{Per-quarter base aggregates (Q1, Q2, and combined H1).} For each
property, the days available, booked, and reserved are counted, with the number of
unique reservations, and a battery of price statistics (mean, median, std, min,
max, 25th/75th percentiles, number of distinct prices). From these are derived
booking/available/reserved rates, price range, coefficient of variation, and
inter-quartile range. Computing them separately for Q1, Q2, and the whole half
year lets the model see both levels and trends.
\item \textbf{Weekend vs.\ weekday split.} Friday--Saturday and Sunday behave
differently from the rest of the week, so each quarter is split into weekend and
weekday subsets and recorded booking rates, mean prices, a weekend price
premium, and the weekend-minus-weekday booking-rate gap.
\item \textbf{Monthly booked days.} The number of booked days in each of the six
months of H1, so the model can pick up acceleration or deceleration heading into
Q3.
\item \textbf{Status transition counts.} How often a listing flips between A, B,
and R, plus counts of A$\to$B (new bookings) and B$\to$A (booking ends). A
listing that changes state often is an active one; a static one barely moves.
\item \textbf{Recency window.} The same kind of aggregates computed only over
the last 30 days of Q2 (June 2016), because the most recent behaviour is the
best clue about what happens next.
\item \textbf{Quarter-over-quarter demand proxy.} The relative change in mean
price from Q1 to Q2 and the change in booking rate, used as a cheap proxy for how
demand was trending.
\item \textbf{Property metadata.} Property age in days at the start of Q3,
reviews per month, price per guest, weekly/monthly discount ratios, a bed/bath
ratio, boolean host flags (Superhost, BusinessReady, InstantbookEnabled), and a
set of missing-indicator columns for the fee/response fields mentioned above.
\item \textbf{Text keywords.} Binary flags are extracted from the listing
title for the presence of words like ``luxury,'' ``private,'' ``near,''
``cozy,'' ``view,'' and ``modern,'' plus the title length and word count.
\item \textbf{Geographic clusters.} K-Means ($k=12$) is run on
latitude/longitude to give every property a neighbourhood cluster label, which
is cheaper and more robust than using the raw zip code as a category.
\end{itemize}

In total this produced 150 features, of which 143 are numeric and 6 are
categorical. Table~\ref{tab:featgroups} summarizes the families and a
representative member of each. The dominant families by count are the
per-quarter behavioural aggregates and the recency window, which is consistent
with the correlation analysis (Fig.~\ref{fig:corr}) and with the
feature-importance ranking reported later (Fig.~\ref{fig:importance}).

\begin{table}[!t]
\caption{Engineered feature families (150 features total: 143 numeric,
6 categorical).}
\label{tab:featgroups}
\centering
\small
\begin{tabular}{lp{1.45in}}
\toprule
Family & Representative feature \\
\midrule
Per-quarter aggregates (Q1/Q2/H1) & reserved-day count, price mean/std \\
Weekend vs.\ weekday & weekend booking-rate gap \\
Monthly booked days & booked days in June 2016 \\
Status transitions & count of A$\to$B flips \\
Recency window (last 30d) & reserved rate, June 2016 \\
QoQ demand proxy & Q1$\to$Q2 booking-rate change \\
Property metadata & property age, reviews/month, host flags \\
Text keywords & ``luxury''/``private'' title flags \\
Geographic & K-Means cluster ($k{=}12$) \\
\bottomrule
\end{tabular}
\end{table}

\begin{figure}[!t]
\centering
\includegraphics[width=\columnwidth]{fig2_top_corr.png}
\caption{Top engineered features by absolute correlation with the target.
Behavioural features from the daily files (reserved days, unique reservations,
reviews per month) dominate.}
\label{fig:corr}
\end{figure}

\subsection{Preprocessing Pipeline and the Target Encoder}
\label{sec:preproc}
Every preprocessing step lives inside a scikit-learn \texttt{ColumnTransformer}
so that it is fit only on the training fold and then applied to the validation
fold, which keeps the workflow leakage-free. Numeric columns are
median-imputed (robust to the price outliers) with missing-indicator
columns where the missing rate is high. Low-cardinality categoricals (15 or
fewer levels, such as \texttt{PropertyType}, \texttt{ListingType},
\texttt{CancellationPolicy}, and the geo cluster) are one-hot encoded.
High-cardinality categoricals (\texttt{Neighborhood},
\texttt{MetropolitanStatisticalArea}) are handled by a custom $K$-fold target
encoder with additive smoothing.

\paragraph{Target-encoding formula} Let $y$ have global mean
$\mu=\tfrac{1}{n}\sum_{i=1}^{n}y_i$. For a categorical column, let level $c$
contain $n_c$ training rows with category mean
$\bar y_c=\tfrac{1}{n_c}\sum_{i:\,x_i=c}y_i$. The smoothed encoding for that
level is the count-weighted average of the category mean and the global prior,
\begin{equation}
e_c \;=\; \frac{n_c\,\bar y_c \;+\; \alpha\,\mu}{\,n_c+\alpha\,},
\qquad \alpha = 20,
\label{eq:smooth}
\end{equation}
where the smoothing factor $\alpha$ is, in effect, a number of ``pseudo-counts''
of the prior: when a level is rare ($n_c\!\ll\!\alpha$) the encoding shrinks
toward $\mu$, and when it is common ($n_c\!\gg\!\alpha$) it approaches the raw
category mean $\bar y_c$. Levels unseen at fit time map to $\mu$.

\paragraph{Leakage-free K-fold scheme} To prevent a row from informing its own
encoding, the encoder is fitted in an internal $K$-fold loop ($K=5$). Writing
$\mathcal{F}_1,\dots,\mathcal{F}_K$ for the inner folds, the encoded value of a
training row $i\in\mathcal{F}_k$ uses statistics computed \emph{only} on the
complementary folds,
\begin{equation}
\tilde e_i \;=\; e^{\,(-\mathcal{F}_k)}_{\,x_i},
\qquad i\in\mathcal{F}_k,
\label{eq:kfoldte}
\end{equation}
where $e^{\,(-\mathcal{F}_k)}_{c}$ is \eqref{eq:smooth} evaluated on
$\{(\mathbf{x}_j,y_j):j\notin\mathcal{F}_k\}$. At transform time (validation and
test rows), the full-data map is used.

To make \eqref{eq:smooth} concrete, consider a neighbourhood that appears in
$n_c=4$ training rows with a local mean of $\bar y_c=40$ booked nights, against a
global mean of $\mu=18.55$. With $\alpha=20$ the encoded value is
$e_c=\frac{4\cdot 40+20\cdot 18.55}{4+20}=\frac{160+371}{24}\approx 22.1$, i.e.\
the estimate is pulled strongly back toward the global prior because only four
listings support it. A neighbourhood with $n_c=400$ rows and the same local mean
would instead encode to $\frac{400\cdot40+20\cdot18.55}{420}\approx 38.9$, almost
the raw mean. This shrinkage is precisely what stops rare high-cardinality levels
from injecting noisy, over-confident estimates into the model. For XGBoost, the model additionally relies on its
native missing-value handling rather than imputing, since it learns the best
split direction for NaNs on its own \cite{chen2016xgboost}. For the linear
baseline, \texttt{StandardScaler} is applied; for the tree models scaling is
unnecessary, but the rest of the pipeline is kept identical so the comparison is
fair.

\subsection{Model Selection and Rationale}
Four model families are compared, and a convex blend is added on top.

\subsubsection{Linear Regression (baseline)}
A plain linear regression provides an interpretable lower bound and indicates which
features are linearly related to the target; its structural limitation on this
bounded target is analysed in Section~\ref{sec:linearcritique}.

\subsubsection{Random Forest Regressor}
A bagging ensemble \cite{breiman1996,breiman2001} that is robust to noise and
needs little tuning. It serves as a strong non-linear reference point.

\subsubsection{Gradient Boosting and XGBoost}
Boosting builds the model sequentially and is widely regarded as the best
off-the-shelf method for tabular data \cite{friedman2001,chen2016xgboost}.
XGBoost minimizes the regularized objective
\begin{equation}
\mathcal{L}(\phi)=\sum_{i=1}^{n}\ell\!\big(y_i,\hat y_i\big)
+\sum_{t}\Omega(f_t),\quad
\Omega(f)=\gamma T+\tfrac{1}{2}\lambda\lVert\mathbf{w}\rVert^{2}
+\alpha\lVert\mathbf{w}\rVert_{1},
\label{eq:xgbobj}
\end{equation}
with squared-error loss $\ell(y,\hat y)=(y-\hat y)^2$, $T$ the number of leaves,
$\mathbf{w}$ the leaf weights, and $\gamma,\lambda,\alpha$ the complexity,
$L_2$, and $L_1$ penalties. These are tuned most aggressively
(Section~\ref{sec:tuning}).

\subsubsection{Multi-Layer Perceptron (MLP)}
\label{sec:mlp}
Finally, a small feed-forward neural network is trained in PyTorch
\cite{paszke2019pytorch} with the Adam optimizer \cite{kingma2015adam}. The network maps the preprocessed input
$\mathbf{z}^{(0)}=\mathbf{x}\in\mathbb{R}^{d}$ through three hidden blocks of
width $(h_1,h_2,h_3)=(256,128,64)$. Each block applies an affine map, batch
normalization, a ReLU non-linearity, and dropout with rate $p=0.2$:
\begin{align}
\mathbf{a}^{(\ell)} &= \mathbf{W}^{(\ell)}\mathbf{z}^{(\ell-1)}+\mathbf{b}^{(\ell)},
\label{eq:mlp_aff}\\
\mathbf{z}^{(\ell)} &= \mathrm{Dropout}_{p}\!\Big(\mathrm{ReLU}\big(
\mathrm{BN}(\mathbf{a}^{(\ell)})\big)\Big),\quad \ell=1,2,3,
\label{eq:mlp_block}
\end{align}
where $\mathrm{ReLU}(u)=\max(0,u)$ and batch normalization standardizes each
unit over the mini-batch $\mathcal{B}$,
\begin{equation}
\mathrm{BN}(a)=\gamma\,\frac{a-\mu_{\mathcal{B}}}{\sqrt{\sigma_{\mathcal{B}}^{2}+\epsilon}}+\beta .
\label{eq:bn}
\end{equation}
A linear read-out produces the scalar prediction in log space,
$\hat r=\mathbf{w}^{(4)\top}\mathbf{z}^{(3)}+b^{(4)}$. Because the count is
skewed, the network is trained on the log-transformed target
$t=\log(1+y)$ with the mean-squared-error criterion
$\mathcal{L}=\tfrac{1}{|\mathcal{B}|}\sum_{i\in\mathcal{B}}(\hat r_i-t_i)^2$, and
predictions are mapped back and clipped,
$\hat y=\mathrm{clip}\big(e^{\hat r}-1,\,0,\,92\big)$. Parameters are updated by
Adam with learning rate $\eta=10^{-3}$ and $L_2$ weight decay $10^{-5}$; using
the standard moment estimates $\mathbf{m}_t,\mathbf{v}_t$ of the gradient
$\mathbf{g}_t$,
\begin{equation}
\boldsymbol{\theta}_{t+1}=\boldsymbol{\theta}_{t}
-\eta\,\frac{\hat{\mathbf m}_t}{\sqrt{\hat{\mathbf v}_t}+\epsilon},
\quad
\hat{\mathbf m}_t=\frac{\mathbf m_t}{1-\beta_1^{t}},\;
\hat{\mathbf v}_t=\frac{\mathbf v_t}{1-\beta_2^{t}}.
\label{eq:adam}
\end{equation}
Training uses a batch size of 512 for up to 60 epochs with early stopping on a
held-out validation MSE (patience 8). Going in, the literature
\cite{tang2024,camatti2024} suggested the trees would win, and they
did---but the MLP turned out to be a useful blend partner because its errors
were not perfectly correlated with the trees' (Section~\ref{sec:complementarity}).

A few design choices deserve a note. Unlike the final XGBoost, the MLP is
trained on the log-transformed target $\log(1+y)$ rather than the raw count.
Neural networks are sensitive to the scale and skew of the target, and the raw
target's long right tail produced unstable gradients and slow convergence; the
log transform compresses that tail so the squared-error loss is better
conditioned, and the monotone inverse $e^{\hat r}-1$ followed by clipping returns
predictions to the count scale. This does not contradict the decision to
\emph{drop} the log transform for XGBoost (Section~\ref{sec:changes}): boosted
trees are invariant to monotone feature scaling and optimize the raw-scale loss
directly, so for them the transform only distorted the objective, whereas for the
MLP it is a numerical-stability aid. Batch normalization
\eqref{eq:bn} keeps the pre-activations well scaled across the three hidden
layers, and dropout ($p=0.2$) plus the $L_2$ weight decay in \eqref{eq:adam}
regularize a network that would otherwise overfit a 19{,}497-row table with 150
inputs.

\subsection{Hyper-Parameter Tuning}
\label{sec:tuning}
All tuning is done with 5-fold cross-validation on the 80\% development split.
For XGBoost, a custom randomized search is used instead of
\texttt{RandomizedSearchCV}, to enable fold-level early stopping: inside
each of the five outer folds, a small internal validation set is carved out, fitting
up to $B_{\max}=3000$ trees with \texttt{early\_stopping\_rounds}~$=50$, and keep
only the best iteration. \texttt{RandomizedSearchCV} does not cleanly allow a
per-fold \texttt{eval\_set} cleanly, so a custom loop was the simplest correct
option, and it also logs every configuration tried.

\paragraph{Search objective} Let $\Theta$ be the hyper-parameter space and let
the data be partitioned into outer folds $\mathcal{D}_1,\dots,\mathcal{D}_K$
($K=5$). For a configuration $\theta$, within fold $k$ the training part is split
into a fitting set and an inner validation set $\mathcal{V}_k$, grow boosting
iterations $b=1,\dots,B_{\max}$, and select the early-stopping iteration
\begin{equation}
b_k^{*}(\theta)=\operatorname*{arg\,min}_{1\le b\le B_{\max}}
\;\mathrm{RMSE}\!\Big(y^{\mathcal{V}_k},\,
\hat f^{(b)}_{\theta,k}\big(\mathbf{X}^{\mathcal{V}_k}\big)\Big),
\label{eq:earlystop}
\end{equation}
stopping once no improvement is seen for 50 consecutive rounds. The search then
minimizes the mean cross-validated MSE of the early-stopped models,
\begin{equation}
\theta^{*}=\operatorname*{arg\,min}_{\theta\in\Theta}\;
\frac{1}{K}\sum_{k=1}^{K}
\mathrm{MSE}\!\Big(y^{\mathcal{D}_k},\,
\hat f^{\,(b_k^{*}(\theta))}_{\theta,k}\big(\mathbf{X}^{\mathcal{D}_k}\big)\Big).
\label{eq:searchobj}
\end{equation}
A total of $N=60$ configurations are drawn from the product of the marginals in
Table~\ref{tab:search}, sampling the learning rate and the regularization terms
log-uniformly so that small and large scales are explored evenly, as recommended
by \cite{bergstra2012random}. The full procedure is summarized in
Algorithm~1.

\begin{table}[!t]
\caption{Search distributions for the 60-configuration randomized search.
$\mathrm{LogU}$ denotes a log-uniform draw, $\mathcal{U}$ a uniform draw.}
\label{tab:search}
\centering
\begin{tabular}{ll}
\toprule
Hyper-parameter & Sampling distribution \\
\midrule
\texttt{n\_estimators} (cap $B_{\max}$) & $3000$ (fixed) \\
\texttt{learning\_rate} & $\mathrm{LogU}(0.02,\,0.10)$ \\
\texttt{max\_depth} & $\mathcal{U}\{4,\dots,8\}$ \\
\texttt{min\_child\_weight} & $\mathcal{U}\{1,\dots,11\}$ \\
\texttt{subsample} & $\mathcal{U}(0.65,\,0.95)$ \\
\texttt{colsample\_bytree} & $\mathcal{U}(0.65,\,0.95)$ \\
\texttt{gamma} ($\gamma$) & $\mathcal{U}(0.0,\,0.4)$ \\
\texttt{reg\_alpha} ($\alpha$) & $\mathrm{LogU}(10^{-3},\,0.5)$ \\
\texttt{reg\_lambda} ($\lambda$) & $\mathrm{LogU}(0.5,\,5.0)$ \\
\bottomrule
\end{tabular}
\end{table}

\vspace{2pt}
\noindent\textbf{Algorithm 1: Randomized search with fold-level early stopping.}
\emph{Input:} hyper-parameter space $\Theta$, outer folds
$\{\mathcal{D}_k\}_{k=1}^{K}$, budget $N$, tree cap $B_{\max}$.
\begin{enumerate}
\item Initialize $\theta^{*}\leftarrow\varnothing$, $s^{*}\leftarrow+\infty$.
\item For $j=1,\dots,N$: sample a configuration $\theta\sim\Theta$ from
Table~\ref{tab:search}.
\item For each fold $k=1,\dots,K$: fit the preprocessing on the training part,
carve an inner validation set $\mathcal{V}_k$, grow up to $B_{\max}$ trees, pick
the early-stop iteration $b_k^{*}$ by \eqref{eq:earlystop}, and record
$m_k=\mathrm{MSE}$ on $\mathcal{D}_k$ using $b_k^{*}$ trees.
\item Set $s=\frac{1}{K}\sum_k m_k$ and log $(\theta,s)$; if $s<s^{*}$, update
$s^{*}\leftarrow s$, $\theta^{*}\leftarrow\theta$.
\item After $N$ configurations, return $\theta^{*}$.
\end{enumerate}
\vspace{2pt}

The best configuration (Table~\ref{tab:xgbparams}) uses a learning rate of
0.027, max depth 6, subsample 0.79, and \texttt{colsample\_bytree} 0.72, and
early stopping settled around 460 trees on average.

\begin{table}[!t]
\caption{Best XGBoost (v3) hyper-parameters found by the 60-configuration
randomized search.}
\label{tab:xgbparams}
\centering
\begin{tabular}{ll}
\toprule
Hyper-parameter & Value \\
\midrule
\texttt{n\_estimators} (cap) & 3000 \\
\texttt{learning\_rate} & 0.0271 \\
\texttt{max\_depth} & 6 \\
\texttt{min\_child\_weight} & 2 \\
\texttt{subsample} & 0.793 \\
\texttt{colsample\_bytree} & 0.718 \\
\texttt{gamma} & 0.268 \\
\texttt{reg\_alpha} & 0.0151 \\
\texttt{reg\_lambda} & 3.401 \\
\texttt{early\_stopping\_rounds} & 50 ($\approx$460 trees kept) \\
\bottomrule
\end{tabular}
\end{table}

\noindent\emph{Note.} In the primary implementation
(\texttt{05\_Gradient\_Boosting\_XGBoost.ipynb}), a 1{,}000-estimator cap is used
for efficient execution. This is consistent with the tuning experiments: despite
a higher search cap of 3{,}000, early stopping consistently identified the
optimal iteration at $\approx$460 trees, confirming that the 1{,}000-estimator
limit is well suited to capture the model's full predictive power without
overfitting.

\subsection{Ensembling (Weighted Blend)}
\label{sec:blend}
Rather than using a single model, the out-of-fold (OOF) predictions are combined
of the candidate models with a convex weighted average. Stack the $m$ models'
OOF predictions into a matrix $\mathbf{P}\in\mathbb{R}^{n\times m}$ with column
$j$ holding model $j$'s OOF prediction. The aim is to find weights
$\mathbf{w}\in\mathbb{R}^{m}$ that minimize the blended MSE, subject to the
predictions staying on the valid scale:
\begin{align}
\min_{\mathbf{w}\in\mathbb{R}^{m}}\quad
& \frac{1}{n}\sum_{i=1}^{n}\Big(y_i-\mathrm{clip}\big[(\mathbf{P}\mathbf{w})_i,\,0,\,92\big]\Big)^{2}
\label{eq:blendobj}\\
\text{s.t.}\quad & \sum_{j=1}^{m}w_j = 1,
\label{eq:blendeq}\\
& 0 \le w_j \le 1,\quad j=1,\dots,m.
\label{eq:blendbox}
\end{align}
The equality constraint~\eqref{eq:blendeq} and the box
bounds~\eqref{eq:blendbox} make the solution a \emph{convex combination}, which
keeps the blend on the same scale as the inputs and prevents any single model
from receiving a runaway weight. Problem \eqref{eq:blendobj}--\eqref{eq:blendbox} is solved
with Sequential Least-Squares Quadratic Programming (SLSQP) via
\texttt{scipy.optimize.minimize}, starting from the uniform point
$w_j^{(0)}=1/m$ with \texttt{ftol}~$=10^{-9}$. As a post-processing step, weights below $10^{-4}$
are zeroed and the rest renormalized so that
\eqref{eq:blendeq} still holds exactly. A randomized search over the
weight simplex yields the same optimum,
which confirms that the SLSQP solution is consistent with the course's core
techniques rather than an artefact of the external optimizer. The optimizer
concentrated most of the weight on XGBoost ($\approx0.73$) and gave smaller,
non-zero weights to the MLP ($\approx0.14$), Linear Regression ($\approx0.10$),
and Random Forest ($\approx0.03$), while Gradient Boosting was dropped to zero
(Table~\ref{tab:weights}). The full procedure is given in Algorithm~2.

\vspace{2pt}
\noindent\textbf{Algorithm 2: Convex weighted blend via SLSQP.}
\emph{Input:} OOF matrix $\mathbf{P}\in\mathbb{R}^{n\times m}$, targets
$\mathbf{y}$.
\begin{enumerate}
\item Define the objective
$g(\mathbf{w})=\frac{1}{n}\lVert\mathbf{y}-\mathrm{clip}(\mathbf{P}\mathbf{w},0,92)\rVert_2^2$.
\item Initialize $\mathbf{w}^{(0)}=(1/m,\dots,1/m)$.
\item Minimize $g$ with SLSQP subject to $\sum_j w_j=1$ and $0\le w_j\le1$ to
obtain $\mathbf{w}^{*}$.
\item Prune: set $w_j^{*}=0$ for every $j$ with $w_j^{*}<10^{-4}$.
\item Renormalize $\mathbf{w}^{*}\leftarrow\mathbf{w}^{*}/\sum_j w_j^{*}$ and
return $\mathbf{w}^{*}$.
\end{enumerate}
\vspace{2pt}

\subsection{Experimental Design and Evaluation}
The labelled data are split once into an 80\% development set (19{,}497
properties) and a 20\% local hold-out (4{,}875 properties) with a fixed seed, so
every model is trained and selected on exactly the same rows. Inside the
development set, 5-fold cross-validation is used for both model selection and
tuning. MSE is optimised directly because it is the evaluation
metric, while MAE and $R^2$ are also tracked internally because they are easier to
read. To prevent leakage, (i) every preprocessing step is fitted only on
the training fold, (ii) no feature is built from Q3 2016 information, and (iii)
target encodings are computed inside their own inner $K$-fold per
\eqref{eq:kfoldte}. As a last step, predictions are clipped to $[0,92]$ and rounded
them to integers, both because the competition requires integer predictions and
because no property can be booked more than 92 nights in a quarter.

\subsection{Computational Cost}
\label{sec:cost}
Because the experiments ran on a single workstation rather than a cluster,
computational cost was monitored throughout. Building the 150-feature matrix from the eight
million daily rows is a one-time aggregation that dominates wall-clock time but
is trivially parallelizable over properties. Model training is cheap by
comparison: a single early-stopped XGBoost fit converges at $\approx$294 trees
on average and takes about 10.6 seconds per 5-fold evaluation on the hardware of Appendix~A, so
the entire 60-configuration randomized search (Algorithm~1)
completed in roughly 10.6 minutes. The Random Forest and the MLP are of the same
order; the linear and gradient-boosting baselines are faster still. The SLSQP
blend itself is essentially free---it optimizes only $m=5$ weights over a fixed
$n\times m$ matrix of cached OOF predictions, converging in well under a second.
This favourable cost profile is one practical reason a small,
well-tuned ensemble is preferred over a single very large model: the marginal compute of
adding the blend is negligible relative to the search, yet it returns a
statistically significant accuracy gain (Section~\ref{sec:significance}).

\subsection{Reproducibility and Implementation}
\label{sec:repro}
Every result in this paper is reproducible from a single fixed seed
(\texttt{random\_state}~$=42$), which governs the 80/20 development/hold-out
split, the 5-fold cross-validation partitions, the inner target-encoding folds,
and the random-search sampler. The pipeline is implemented entirely in Python
with scikit-learn \cite{pedregosa2011sklearn} for the preprocessing
\texttt{ColumnTransformer}, imputers, one-hot encoder, linear, random-forest and
gradient-boosting models and cross-validation; the \texttt{xgboost} package for
the boosted trees \cite{chen2016xgboost}; PyTorch \cite{paszke2019pytorch} with
Adam \cite{kingma2015adam} for the MLP; and \texttt{scipy.optimize} for the SLSQP
blend. The code is organized as a sequence of numbered notebooks---exploratory
analysis, feature engineering, one notebook per model family, and the blend---so
that each stage can be re-run independently, and the out-of-fold predictions are
cached to disk (\texttt{.npy}) so that the blend and all of the post-hoc
diagnostics in Section~\ref{sec:results} can be recomputed without retraining.

%==============================================================================
\section{Results and Discussion}
\label{sec:results}
Table~\ref{tab:cv} lists the 5-fold cross-validation scores of all models on the
development set, including two strong external gradient-boosting baselines. The
ordering is consistent with the literature. Linear Regression and the
lightly-tuned Gradient Boosting trail the field, Random Forest and the MLP sit in
the middle, and the boosted-tree family is strongest. After the 60-configuration
retuning with early stopping, the tuned XGBoost regressor reaches a 5-fold CV MSE
of 182.2, and the SLSQP blend lowers the out-of-fold error to 180.0---about a
$2\%$ relative improvement over the best single model.

Crucially, the proposed pipeline is also benchmarked against two
state-of-the-art gradient-boosting libraries widely used in tabular forecasting,
LightGBM and CatBoost, trained with the same leakage-safe preprocessing and the
same 5-fold split. As reported in Table~\ref{tab:cv}, tuned LightGBM (184.1) and
CatBoost (184.3) are competitive but do \emph{not} surpass the in-house tuned
XGBoost (182.2), and both are clearly behind the convex blend (180.0). This
shows that the reported gains stem from the temporal feature engineering, the
leakage-safe target encoding, and the optimized convex ensemble rather than from
the choice of a particular boosting library.

\begin{table}[!t]
\caption{Model comparison by 5-fold CV on the development set (19{,}497
properties; lower MSE/MAE and higher $R^2$ are better). LightGBM and CatBoost are
state-of-the-art external baselines trained with the identical pipeline.}
\label{tab:cv}
\centering
\begin{tabular}{lccc}
\toprule
Model & MSE & MAE & $R^2$ \\
\midrule
LinearRegression & 207.6 & 9.22 & 0.731 \\
GradientBoosting & 245.3 & 8.46 & 0.683 \\
RandomForest & 194.2 & 8.55 & 0.749 \\
MLP & 217.6 & 8.03 & 0.718 \\
\midrule
LightGBM (SOTA baseline) & 184.1 & 8.24 & 0.762 \\
CatBoost (SOTA baseline) & 184.3 & 8.27 & 0.762 \\
XGBoost (tuned, proposed) & 182.2 & 8.18 & 0.764 \\
\textbf{SLSQP weighted blend (proposed)} & \textbf{180.0} & \textbf{8.11} & \textbf{0.767} \\
\bottomrule
\end{tabular}
\end{table}

\begin{figure}[!t]
\centering
\includegraphics[width=\columnwidth]{fig3_baseline_cv.png}
\caption{Baseline model comparison by cross-validated MSE. XGBoost is the
strongest single family before any dedicated retuning.}
\label{fig:baseline}
\end{figure}

\subsection{Why the Methods Behaved as They Did}
The gap between the linear model and the boosted trees is the clearest result,
and it is exactly the structural mismatch made precise in
Section~\ref{sec:linearcritique}: the relationship between the features and the
bounded, zero-inflated target is strongly non-linear and full of
interactions---for example, a high reserved rate in Q2 means something very
different for a brand-new listing than for one with hundreds of reviews. Linear
regression cannot represent those interactions without hand-crafting them, while
trees discover them automatically. Random Forest captures the non-linearity but,
being a variance-reduction (bagging) method, it cannot reduce bias as
aggressively as boosting can. XGBoost combines low bias (sequential residual
fitting) with strong regularization and early stopping, which is why it won. The
MLP landed between the trees and the linear model; neural networks usually need
more data and more careful tuning to beat gradient boosting on tabular problems,
and the compute budget did not allow a deep search.

It is useful to read the whole ranking through a bias--variance lens. Linear
regression is a high-bias, low-variance estimator: it cannot bend to the
saturating conditional mean of Section~\ref{sec:linearcritique}, so its error is
dominated by bias and no amount of data removes it. Random Forest attacks the
opposite term---bagging averages many high-variance trees to cut variance---but,
because every tree is grown to near-purity on the same features, it leaves a
residual bias that boosting can still reduce. Gradient boosting lowers bias by
fitting residuals stage by stage, and XGBoost adds the regularization and early
stopping \eqref{eq:xgbobj}--\eqref{eq:earlystop} that keep the corresponding
variance in check, which is why it sits at the bottom of the MSE column. The MLP
is capable of low bias in principle, but on 19{,}497 rows it is variance-limited
without heavier regularization, which is exactly why batch normalization and
dropout (Section~\ref{sec:mlp}) mattered so much to its stability. The blend then
performs one more variance reduction \emph{across} model families rather than
within one, which is why it improves on even the best single model.

\begin{figure}[!t]
\centering
\includegraphics[width=\columnwidth]{fig4_xgb_importance.png}
\caption{XGBoost gain-based feature importances. Behavioural and recency
features dominate, in agreement with Fig.~\ref{fig:corr}.}
\label{fig:importance}
\end{figure}

The feature-importance analysis (Fig.~\ref{fig:importance}) confirms the EDA
story: the model leans most heavily on the Q2 reserved-day and reservation
counts, reviews per month, and the recency-window features, with the static
metadata and text keywords playing only a supporting role. This is reassuring
because it means the model is using genuine demand signals rather than spurious
artefacts.

\subsection{Impact of the Blend and the Optimization Choices}
Two optimization decisions mattered most for the final score. The first was
early stopping: capping the number of trees at 3{,}000 but letting each fold stop
when its internal validation error stopped improving captured most of the benefit
of a large ensemble without the overfitting, and it made the search cheaper
because bad configurations stopped early. The second was the blend. An
equal-weight average of the five models scored 186.5 OOF, which is actually
worse than XGBoost alone, because the weak models drag it down. Once
SLSQP chooses the weights, the error drops to 180.0---the optimizer effectively
ignored the weak models and kept only the complementary ones.
Fig.~\ref{fig:blend} shows the per-model and blended OOF scores side by side, and
Table~\ref{tab:weights} lists the optimal weights.

\begin{figure}[!t]
\centering
\includegraphics[width=\columnwidth]{fig5_blend.png}
\caption{Out-of-fold MSE for each model and for the equal-weight vs.\ optimal
(SLSQP) blends. The weighted blend is the best configuration.}
\label{fig:blend}
\end{figure}

\begin{table}[!t]
\caption{Optimal blend weights chosen by SLSQP on the OOF predictions
(problem~\eqref{eq:blendobj}--\eqref{eq:blendbox}).}
\label{tab:weights}
\centering
\begin{tabular}{lc}
\toprule
Model & Weight $w_j$ \\
\midrule
LinearRegression & 0.0628 \\
RandomForest & 0.0000 \\
GradientBoosting & 0.0000 \\
XGBoost & 0.3406 \\
MLP & 0.1064 \\
XGBoost v3 & 0.2557 \\
LightGBM & 0.2014 \\
CatBoost & 0.0330 \\
\bottomrule
\end{tabular}
\end{table}

\subsection{Model Complementarity and Final Error Metrics}
\label{sec:complementarity}
A blend can only help if its members make \emph{different} mistakes; if every
model erred on the same rows, averaging them would just reproduce the common
error. Table~\ref{tab:errcorr} reports the Pearson correlation of the per-row
prediction errors across the five models. The tree models are highly correlated
with one another (XGBoost--Random Forest reaches $0.96$), which is why the
optimizer keeps only one of them at full weight and drives Gradient Boosting to
zero (Table~\ref{tab:weights}). The MLP and Linear Regression are the least
correlated with XGBoost ($0.89$ and $0.91$), so although they are individually
weaker (Table~\ref{tab:metrics}), they carry independent signal and receive the
non-zero supporting weights. This is the quantitative justification for the
weight pattern in Table~\ref{tab:weights}: SLSQP is not merely picking the best
model, it is exploiting error de-correlation.

Table~\ref{tab:metrics} also reports the blend's MAE and $R^2$ alongside MSE.
The blend improves not only the optimized MSE (180.01 vs.\ 182.19) but also the
$R^2$ (0.766 vs.\ 0.764), and its MAE (8.10) is within a hair of the best
individual MAE (the MLP's 8.03)---so the MSE gain does not come at the cost of
median accuracy.

\begin{table}[!t]
\caption{Final error metrics on the development OOF predictions (lower MSE/MAE,
higher $R^2$ are better).}
\label{tab:metrics}
\centering
\begin{tabular}{lccc}
\toprule
Model & MSE & MAE & $R^2$ \\
\midrule
LinearRegression & 207.61 & 9.22 & 0.731 \\
RandomForest & 194.24 & 8.55 & 0.749 \\
GradientBoosting & 245.29 & 8.46 & 0.683 \\
MLP & 220.26 & 8.03 & 0.718 \\
XGBoost & 182.19 & 8.18 & 0.764 \\
XGBoost v3 & 182.87 & 8.18 & 0.763 \\
LightGBM & 184.39 & 8.24 & 0.762 \\
CatBoost & 184.55 & 8.27 & 0.762 \\
\textbf{SLSQP blend} & \textbf{180.01} & \textbf{8.11} & \textbf{0.767} \\
\bottomrule
\end{tabular}
\end{table}

\begin{table}[!t]
\caption{Pearson correlation of the per-row \emph{prediction errors} across
models (development OOF). Lower off-diagonal values mark more complementary
models; the matrix is symmetric so only the upper triangle is shown.}
\label{tab:errcorr}
\centering
\begin{tabular}{lccccc}
\toprule
 & LR & RF & GB & XGB & MLP \\
\midrule
LR  & 1.00 & 0.91 & 0.84 & 0.91 & 0.86 \\
RF  &      & 1.00 & 0.88 & 0.96 & 0.87 \\
GB  &      &      & 1.00 & 0.89 & 0.87 \\
XGB &      &      &      & 1.00 & 0.89 \\
MLP &      &      &      &      & 1.00 \\
\bottomrule
\end{tabular}
\end{table}

\subsection{Is the Blend Gain Statistically Significant?}
\label{sec:significance}
A drop from 182.2 to 180.0 MSE is small, so it is tested whether the gain is real or just
cross-validation noise. Because both models produce predictions for the same
19{,}497 development rows, a \emph{paired} test on the per-sample squared
errors is used, the natural sample-level analogue of the paired resampling tests
recommended for model comparison \cite{dietterich1998}. Let
$e_i^{\text{xgb}}=y_i-\hat y_i^{\text{xgb}}$ and
$e_i^{\text{bl}}=y_i-\hat y_i^{\text{bl}}$, and define the per-row difference of
squared errors
\begin{equation}
d_i=\big(e_i^{\text{xgb}}\big)^{2}-\big(e_i^{\text{bl}}\big)^{2}.
\label{eq:diff}
\end{equation}
Its sample mean $\bar d$ is exactly the MSE gap, $\bar d = 182.19-180.01=2.18$.
A two-sided paired $t$-test of $H_0:\mathbb{E}[d_i]=0$ gives $t=5.33$ with
$p=1.0\times10^{-7}$, and the 95\% confidence interval on the mean squared-error
reduction is $[1.38,\,2.98]$, which excludes zero. The non-parametric Wilcoxon
signed-rank test agrees overwhelmingly ($p<10^{-21}$). $H_0$ is therefore rejected:
the blend's improvement over XGBoost is statistically significant, not an
artefact of CV variance. The effect size is small (Cohen's $d_z\approx0.04$),
which is expected---the blend changes only the minority of rows where the
component models disagree---but it is consistent and free, since the components
were already trained. For full rigour a $5\times2$ cross-validation paired
$t$-test \cite{dietterich1998} could be run by repeating the whole pipeline over
five $2$-fold splits; the sample-level paired test above tests the same
hypothesis on the existing held-out predictions and is what the compute budget
allowed.

\subsection{Residual Analysis on the [0, 92] Boundaries}
\label{sec:residuals}
Because the target is bounded and zero-inflated, aggregate MSE can hide
systematic behaviour at the edges, so the residuals are examined
$r_i=y_i-\hat y_i$ of the blend directly. Fig.~\ref{fig:residuals} shows the
standard diagnostics. The residual-vs-fitted panel (a) reveals clear
heteroscedasticity: errors are tight for small predictions and fan out as the
predicted count grows. The residual-vs-true panel (b) is the most telling for a
censored target: residuals are mildly positive across the low range and become
strongly positive in the upper band $[90,92]$, meaning the model
\emph{under-predicts} the busiest listings. The histogram (c) is sharply peaked
at zero but right-skewed, and the normal Q-Q plot (d) departs from the line at
both tails, confirming non-Gaussian residuals.

\begin{figure}[!t]
\centering
\includegraphics[width=\columnwidth]{fig6_residuals.png}
\caption{Residual diagnostics for the blend on the development OOF predictions:
(a) residual vs.\ fitted, (b) residual vs.\ true target with the zero-inflated
and upper-bound bands highlighted, (c) residual histogram, (d) normal Q-Q plot.}
\label{fig:residuals}
\end{figure}

To localize the bias, the rows are binned by their true value and the mean
signed residual per bin (Fig.~\ref{fig:residualbins}) is plotted. The pattern is exactly
what censoring predicts: for the dead and near-dead listings the model
\emph{over-predicts} (mean residuals of $-3.5$ at $y=0$ and $-8.3$ on $[1,5]$,
i.e.\ $\hat y>y$), while for the busy listings it \emph{under-predicts}
increasingly toward the bound (mean residual $+12.0$ on $[51,80]$, $+15.2$ on
$[81,89]$, and $+23.7$ on $[90,92]$). Clipping to $[0,92]$ removes impossible
predictions but does nothing about this regression-to-the-mean at the
boundaries---the estimator shrinks extreme counts toward the centre of the
distribution. This is the single clearest motivation for the two-stage hurdle
model proposed in Section~\ref{sec:concl}: a dedicated classifier for ``will
this property be booked at all?'' would directly target the over-prediction on
the large zero mass, and a separate regressor on the positive counts could be
trained to be less biased near the upper bound.

\begin{figure}[!t]
\centering
\includegraphics[width=\columnwidth]{fig7_residual_bins.png}
\caption{Mean signed residual ($y-\hat y$) by true booked-nights bin. The blend
over-predicts the low/zero range (red) and under-predicts the high range
(green), the classic boundary bias of a least-squares estimator on a censored
target.}
\label{fig:residualbins}
\end{figure}

\subsection{Effect of Clipping and Integer Rounding}
\label{sec:rounding}
The competition expects integer predictions, and no listing can exceed 92 booked
nights, so the final post-processing step clips each prediction to $[0,92]$ and
rounds it to the nearest integer. It is worth checking that this step does not
quietly cost accuracy. On the development OOF predictions, the floating-point
blend scores 180.010 MSE; after clipping and integer rounding the score rises only
to 180.074, an increase of $+0.073$ MSE, or about $0.04\%$. The penalty is
negligible because the blend's predictions are already close to integers for the
large zero-and-low mass, and clipping only touches the handful of predictions that
stray outside $[0,92]$. This $0.07$-MSE rounding cost, together with the
held-out test set being a different sample of properties, accounts for the small
gap between the internal OOF estimate (180.0) and the held-out test MSE (185.5)
discussed next; reassuringly, the rounding introduces no large or systematic
error.

\subsection{Held-Out Test Performance}
The integer-rounded, clipped predictions of the weighted blend are evaluated on
the held-out test set, where they obtain an MSE of 185.499. This held-out MSE
(185.5) is very close to the internal OOF estimate (180.0); the small gap is
expected, since the test set is a different sample of properties and integer
rounding adds a little error. The close agreement between the two numbers
indicates that the cross-validation procedure is honest and that the model does
not overfit the evaluation set. For reference, a single tuned XGBoost regressor
without any blending obtains 187.740 on the same held-out set, so the convex
blend improves the held-out MSE by roughly $1.2\%$ on truly unseen data, mirroring
the out-of-fold gain.

\subsection{Failure-Case Analysis}
\label{sec:failure}
Aggregate metrics and binned residuals describe the model's behaviour
statistically; a sample-level dissection of its worst errors gives complementary,
actionable insight. Table~\ref{tab:failure} lists the development-set listings
with the largest blend residuals. Every one of them is a property that was booked
for the entire quarter ($y=92$) yet predicted at fewer than seven nights. Their
common signature is a \emph{decoupling} of Q3 demand from the Q1--Q2 behavioural
history on which the model relies: listing~5117519 has a high Q2 reserved rate
($0.60$) but only a single review, listing~1260093 is an expensive ($\$465$/night)
well-reviewed apartment whose Q2 reserved rate is nevertheless only $0.12$, and
listing~12172839 has essentially no review history at all. In each case a sudden
Q3 surge in demand is not foreshadowed by the recent behavioural features, so a
purely historical regressor cannot anticipate it. Aggregated over the $211$
listings with $y\ge90$, the blend predicts a mean of only $67.4$ nights despite a
mean Q2 reserved rate of $0.86$, confirming the systematic under-prediction near
the upper bound. These failures point directly to the value of the two-stage
hurdle model proposed in Section~\ref{sec:concl} and to richer
demand-acceleration features.

\begin{table}[!t]
\caption{Largest blend residuals on the development set. All are fully-booked
listings ($y=92$) that the model severely under-predicts; their recent
behavioural features do not foreshadow the Q3 surge.}
\label{tab:failure}
\centering
\small
\begin{tabular}{rrrrr}
\toprule
Property & $y$ & $\hat y$ & Q2 res.\ rate & \#Reviews \\
\midrule
5117519  & 92 & 2.2 & 0.60 & 1 \\
1260093  & 92 & 3.9 & 0.12 & 68 \\
1625045  & 92 & 5.9 & 0.00 & 17 \\
12172839 & 92 & 6.3 & 0.09 & 0 \\
\bottomrule
\end{tabular}
\end{table}

\subsection{Threats to Validity}
\label{sec:threats}
This subsection states the main threats to validity. \emph{Internal validity:}
all preprocessing is fitted inside the training fold and target encodings use an
inner $K$-fold (Section~\ref{sec:preproc}), preventing train/validation leakage;
the close agreement between the OOF MSE (180.0) and the held-out test MSE (185.5)
is empirical support for this, and integer rounding accounts for part of the
small remaining gap. \emph{External validity:} the model is fitted to a single
market (New York City) and a single quarter (Q3 2016), so the learned price and
seasonality effects need not transfer to other cities or years, and the reliance
on Q1--Q2 behavioural history degrades performance for brand-new listings with no
daily record (Section~\ref{sec:failure}). \emph{Construct validity:} the
significance test in Section~\ref{sec:significance} is a paired test on the
shared OOF predictions rather than a full $5\times2$ cross-validation
\cite{dietterich1998}; it tests the right hypothesis but does not capture
variance from re-drawing the folds themselves.

%==============================================================================
\section{Ablation Study and Design Choices}
\label{sec:changes}
This section isolates the contribution of the proposed components and reports the
key design decisions that shaped the final pipeline.
Table~\ref{tab:ablation} reports a cumulative ablation in which pipeline
components are added one at a time, all measured under the identical 5-fold
cross-validation protocol; the paragraphs that follow analyse the remaining
design choices.

\begin{table}[!t]
\caption{Cumulative ablation (5-fold CV MSE, identical split). The upper block
varies the feature/preprocessing axis with a fixed linear model; the lower block
varies the model/ensemble axis under the full preprocessing.}
\label{tab:ablation}
\centering
\begin{tabular}{lcc}
\toprule
Configuration & MSE & $\Delta$ \\
\midrule
Linear, raw static features only & 346.6 & --- \\
\quad $+$ engineered temporal features & 207.3 & $-139.3$ \\
\quad $+$ K-fold target encoding (full prep.) & 207.4 & $+0.1$ \\
\midrule
XGBoost (tuned, full prep.) & 182.2 & $-25.2$ \\
Naive equal-weight blend (5 learners) & 186.5 & $+4.3$ \\
\textbf{SLSQP convex blend (proposed)} & \textbf{180.0} & $-6.0$ \\
\bottomrule
\end{tabular}
\end{table}

\paragraph{Engineered temporal features dominate} The ablation in
Table~\ref{tab:ablation} is unambiguous about where the accuracy comes from.
Adding the engineered temporal aggregates to a linear model trained on the raw
static property fields alone cuts the MSE from 346.6 to 207.3---a $40\%$
reduction and by far the largest single improvement in the pipeline. This
confirms that the daily listing histories, not the static metadata, carry most of
the predictive signal.

\paragraph{Target encoding is leakage-safety, not accuracy, for the linear model}
Adding the leakage-safe $K$-fold target encoding changes the linear model's MSE
only marginally ($207.3\!\rightarrow\!207.4$). Its value is therefore not a direct
accuracy gain for the linear baseline but the ability to represent
high-cardinality location fields \emph{without leakage}; this matters far more
for the tree models, which exploit those encoded splits, than for the linear one.

\paragraph{Model choice and convex weighting} Switching from the linear model to
the tuned XGBoost under the same preprocessing lowers the MSE by a further $25.2$
(207.4 to 182.2). The ensemble axis confirms that naive equal-weight averaging
(186.5) is \emph{worse} than the best single model, because the weaker learners
drag the mean down; only the constrained convex optimisation of
\eqref{eq:blendobj}--\eqref{eq:blendbox} produces a net gain (180.0), by
concentrating weight on the complementary models and assigning zero weight to the
redundant Gradient Boosting learner (Table~\ref{tab:weights}). Clipping and
integer rounding add a negligible $+0.07$ MSE.

\paragraph{Raw target with clipping beats a log-transform} Training the boosted
trees on a log-transformed target $\log(1+y)$ and inverting at prediction time
optimises the wrong objective, because the evaluation metric is MSE on the raw
count scale while the target is bounded and zero-inflated (Fig.~\ref{fig:target}).
Squared error in log space over-weights the small counts and under-weights the
large counts that dominate the raw MSE. The final XGBoost is therefore trained
directly on the raw target with squared-error loss, and the skew is controlled by
clipping predictions to $[0,92]$.

\paragraph{Keyword flags suffice; TF-IDF does not help} A TF-IDF representation of
the listing titles was evaluated and did not improve accuracy
\cite{islam2022}. Titles are very short, so TF-IDF produces a wide, sparse matrix
with little signal beyond the simple keyword flags, while increasing the risk of
overfitting and the training time. The compact keyword flags
(luxury/private/near/cozy/view/modern) plus title length and word count are
retained, and full text modelling is excluded.

\paragraph{Host-level encoding is omitted} Target or leave-one-out encoding of a
host identifier is a classic leakage and overfitting trap: the identifier has
extremely high cardinality, many hosts appear only once or twice, and the encoded
value would effectively memorise individual hosts. The location and host signal
is captured more safely through the neighbourhood and metropolitan-area target
encodings and the geographic-cluster feature, so the host identifier is excluded.

\paragraph{Custom search enables fold-level early stopping} A hand-written
randomized search is used instead of an off-the-shelf cross-validated search
because fold-level early stopping with an XGBoost \texttt{eval\_set} cannot be
expressed cleanly by the latter. The custom loop preserves the random-sampling
philosophy of \cite{bergstra2012random} and \eqref{eq:searchobj} but additionally
lets each fold stop at its own best iteration \eqref{eq:earlystop} and logs every
trial for inspection.

%==============================================================================
\section{Conclusions}
\label{sec:concl}
This work presented a complete, leakage-safe machine-learning pipeline for
forecasting quarterly short-term-rental demand in New York City, a censored,
zero-inflated regression task. From five raw tables, 150 per-property features
were engineered; five regressors were compared under 5-fold cross-validation;
XGBoost was tuned with an early-stopping randomized search; and the learners were
combined through an SLSQP-optimised convex blend. The blend attains an
out-of-fold MSE of 180.0 and a held-out test MSE of 185.499, outperforming the
best single model and two state-of-the-art gradient-boosting baselines
(LightGBM and CatBoost); the improvement is statistically significant
(Section~\ref{sec:significance}).

Three findings carry beyond this dataset. Most of the predictive power comes from
the daily listing histories rather than the static property table, so temporal
aggregation and recency features are far more valuable than static metadata.
Matching the loss to the evaluation metric (raw MSE with clipping rather than a
log-transform) matters as much as the choice of learner on a bounded,
zero-inflated target. Finally, a constrained convex blend of complementary models
is a low-cost, low-risk way to obtain a statistically significant accuracy gain.

\emph{Scientific and industrial impact.} Accurate quarter-ahead demand
forecasts have concrete value beyond a single benchmark. Platforms can use
booking-density predictions to balance customer-support and infrastructure
capacity ahead of demand peaks, to target marketing toward listings predicted to
be under-booked, and to drive dynamic pricing. At a city scale, booking-density
models can support municipal studies of housing-market pressure and the
displacement of long-term rental stock by short-term lets. The methodological
contribution---formalised leakage-safe target encoding combined with a
constrained convex ensemble for a censored, zero-inflated target---transfers to
other bounded count-forecasting problems such as occupancy, attendance, or
inventory demand.

\subsection{Limitations and Future Work}
\label{sec:limitations}
Several limitations bound the present results and motivate future work.
\emph{Geographic and temporal specificity:} the model is fitted to a single
market (New York City) and a single quarter (Q3 2016), so the learned price and
seasonality effects need not transfer to other cities or years.
\emph{Cold-start dependence:} predictions rely heavily on Q1--Q2 behavioural
history, so accuracy degrades for brand-new listings, which is the dominant
failure mode identified in Section~\ref{sec:failure}.
\emph{Search cost:} the customised fold-level randomized search is inherently
more sequential than a fully parallelised grid search, trading wall-clock
efficiency for the ability to apply per-fold early stopping.
\emph{Static quarterly horizon:} the formulation predicts a single quarterly
aggregate and does not capture dynamic, week-to-week variation in booking
behaviour.

Future work follows directly from these limitations and from the residual
analysis (Section~\ref{sec:residuals}). First, a two-stage hurdle model---a
classifier for whether a property is booked at all, followed by a regressor for
the positive count---would address the zero-inflation and boundary bias more
directly than clipping. Such a model factorises the prediction into a Bernoulli
hurdle and a positive-count expert,
\begin{equation}
\mathbb{E}[y\mid\mathbf{x}] = \underbrace{\Pr(y>0\mid\mathbf{x})}_{\text{classifier }p(\mathbf{x})}\;\cdot\;\underbrace{\mathbb{E}[y\mid y>0,\mathbf{x}]}_{\text{regressor }r(\mathbf{x})},
\label{eq:hurdle}
\end{equation}
where $p(\mathbf{x})$ is trained on the binary label $\mathbb{1}[y>0]$ over all
rows---directly attacking the $53.5\%$ zero mass---and $r(\mathbf{x})$ is fit only
on booked listings, so it is not biased by the zeros that cause the
over-prediction observed on the $[1,5]$ bin (Fig.~\ref{fig:residualbins}).
Second, richer recency features---for example the slope of the booking trend over
the final weeks of Q2---would distinguish listings whose demand is
\emph{accelerating} into Q3 from those merely busy on average, directly targeting
the failure cases in Section~\ref{sec:failure}. Third, because the MLP is the
learner least correlated with XGBoost (Table~\ref{tab:errcorr}), a more careful
neural-network search is the most promising route to a stronger blend partner.

%==============================================================================
\bibliographystyle{IEEEtran}
\bibliography{references}

%==============================================================================
\appendices
\section{Reproducibility and Execution Environment}
\label{app:env}
All results are reproducible from a single fixed seed
(\texttt{random\_state}~$=42$), which governs the 80/20 development/hold-out
split, the 5-fold cross-validation partitions, the inner target-encoding folds,
and the random-search sampler. The pipeline is implemented in Python~3.13 with
scikit-learn~1.8.0 (preprocessing, linear, random-forest and gradient-boosting
models, cross-validation), XGBoost~3.2.0, PyTorch~2.12 (CPU build) with Adam for
the MLP, SciPy~1.17.0 for the SLSQP blend, and LightGBM~4.6.0 and CatBoost~1.2.10
for the external baselines. All experiments were run on a single workstation with
an AMD (Zen~3) multi-core CPU and no GPU. Representative wall-clock times are:
one 5-fold XGBoost evaluation $\approx$10.6\,s, the full 60-configuration search
$\approx$10.6\,min, a 5-fold LightGBM evaluation $\approx$18\,s, a 5-fold CatBoost
evaluation $\approx$40\,s, and the SLSQP blend $<$1\,s. Out-of-fold predictions
are cached to disk so that the blend and all post-hoc diagnostics can be
recomputed without retraining.

\section{Engineered Feature Inventory}
\label{app:features}
Table~\ref{tab:featinventory} documents the 150 engineered features (143 numeric,
7 categorical) by family, with their source tables and a representative
mathematical definition, so that the feature matrix can be reproduced. Per-quarter
aggregates are computed independently for Q1, Q2, and the combined first half
(H1); rates are normalised by the number of days a listing is observed in the
period.

\begin{table}[!t]
\caption{Engineered feature families, source tables, and representative
definitions. ``daily'' denotes the Q1/Q2 daily listing tables; ``static''
denotes the property table. $\mathbb{1}[\cdot]$ is the indicator function.}
\label{tab:featinventory}
\centering
\footnotesize
\begin{tabular}{p{1.05in}p{0.45in}p{1.25in}}
\toprule
Family (source) & \# & Representative definition \\
\midrule
Per-quarter aggregates (daily) & $\sim$85 & booking rate $=\frac{\#\{\text{booked days}\}}{\#\{\text{observed days}\}}$; price CV $=\sigma_{p}/\mu_{p}$ \\
Weekend/weekday split (daily) & 4 & weekend premium $=\mu_{p}^{\text{wknd}}-\mu_{p}^{\text{wkdy}}$ \\
Monthly booked days (daily) & 9 & $\sum_{d\in\text{month}}\mathbb{1}[\text{status}_d=B]$ \\
Status transitions (daily) & 3 & $\#\{d:\text{status}_d{=}A,\text{status}_{d+1}{=}B\}$ \\
Recency window (daily) & 7 & aggregates over last 30 days of Q2 \\
QoQ demand proxy (daily) & 2 & $\frac{\mu_{p}^{Q2}-\mu_{p}^{Q1}}{\mu_{p}^{Q1}}$; $\Delta$ booking rate \\
Property metadata (static) & $\sim$25 & age in days; reviews/month; bed/bath ratio; missing-indicators \\
Text keywords (static) & 10 & $\mathbb{1}[\text{title contains ``luxury''}]$, etc.; length; word count \\
Geographic (static) & 4 & K-Means$(k{=}12)$ cluster of (lat, lon) \\
\bottomrule
\end{tabular}
\end{table}


\section{Complete Feature Inventory}
\label{app:allfeatures_table}
Here is the complete inventory of all 149 engineered features, along with the identifier and target variables, and their corresponding data types.

\begin{table*}[!ht]
\caption{Detailed list of all 151 columns in the processed feature matrix.}
\label{tab:allfeatures_list}
\centering
\footnotesize
\begin{tabular}{ll|ll|ll}
\toprule
Feature Name & Type & Feature Name & Type & Feature Name & Type \\
\midrule
PropertyID & integer (ID) & q1\_days\_A & binary & h1\_available\_rate & float \\
NumReserveDays2016Q3 & integer (target) & q1\_days\_B & binary & h1\_reserved\_rate & float \\
PropertyType & str & q1\_days\_R & binary & h1\_price\_range & float \\
ListingType & str & q1\_unique\_reservations & integer & h1\_price\_cv & float \\
Zipcode & integer & q1\_price\_mean & float & h1\_price\_iqr & float \\
Neighborhood & str & q1\_price\_median & float & q1\_we\_days & integer \\
MetropolitanStatisticalArea & str & q1\_price\_std & float & q1\_we\_booked & binary \\
NumberofReviews & float & q1\_price\_nunique & integer & q1\_we\_price\_mean & float \\
OverallRating & float & q1\_price\_min & float & q1\_wd\_days & integer \\
Bedrooms & float & q1\_price\_max & float & q1\_wd\_booked & binary \\
Bathrooms & float & q1\_price\_q25 & float & q1\_wd\_price\_mean & float \\
MaxGuests & integer & q1\_price\_q75 & float & q1\_we\_booking\_rate & float \\
ResponseRate & float & q1\_booking\_rate & float & q1\_wd\_booking\_rate & float \\
ResponseTimemin & float & q1\_available\_rate & float & q1\_weekend\_premium & float \\
Superhost & float & q1\_reserved\_rate & float & q1\_weekend\_minus\_weekday\_book\_rate & float \\
CancellationPolicy & str & q1\_price\_range & float & q2\_we\_days & integer \\
SecurityDeposit & float & q1\_price\_cv & float & q2\_we\_booked & binary \\
CleaningFee & float & q1\_price\_iqr & float & q2\_we\_price\_mean & float \\
ExtraPeopleFee & float & q2\_days\_total & integer & q2\_wd\_days & integer \\
PublishedNightlyRate & integer & q2\_days\_A & binary & q2\_wd\_booked & binary \\
PublishedMonthlyRate & float & q2\_days\_B & binary & q2\_wd\_price\_mean & float \\
PublishedWeeklyRate & float & q2\_days\_R & binary & q2\_we\_booking\_rate & float \\
MinimumStay & integer & q2\_unique\_reservations & integer & q2\_wd\_booking\_rate & float \\
NumberofPhotos & float & q2\_price\_mean & float & q2\_weekend\_premium & float \\
BusinessReady & float & q2\_price\_median & float & q2\_weekend\_minus\_weekday\_book\_rate & float \\
InstantbookEnabled & float & q2\_price\_std & float & q1\_booked\_m1 & binary \\
Latitude & float & q2\_price\_nunique & integer & q1\_booked\_m2 & binary \\
Longitude & float & q2\_price\_min & float & q1\_booked\_m3 & binary \\
PropertyAgeDays & integer & q2\_price\_max & float & q2\_booked\_m4 & binary \\
ReviewsPerMonth & float & q2\_price\_q25 & float & q2\_booked\_m5 & binary \\
PricePerGuest & float & q2\_price\_q75 & float & q2\_booked\_m6 & binary \\
WeeklyDiscount & float & q2\_booking\_rate & float & q1\_status\_changes & integer \\
MonthlyDiscount & float & q2\_available\_rate & float & q1\_new\_bookings & integer \\
BedBathRatio & float & q2\_reserved\_rate & float & q1\_booking\_ends & integer \\
ResponseRate\_missing & binary & q2\_price\_range & float & q2\_status\_changes & integer \\
ResponseTimemin\_missing & binary & q2\_price\_cv & float & q2\_new\_bookings & integer \\
SecurityDeposit\_missing & binary & q2\_price\_iqr & float & q2\_booking\_ends & integer \\
CleaningFee\_missing & binary & h1\_days\_total & integer & h1\_status\_changes & integer \\
ExtraPeopleFee\_missing & binary & h1\_days\_A & binary & h1\_new\_bookings & integer \\
OverallRating\_missing & binary & h1\_days\_B & binary & h1\_booking\_ends & integer \\
HasReviews & binary & h1\_days\_R & binary & recent\_days & float \\
kw\_luxury & binary & h1\_unique\_reservations & integer & recent\_booked & binary \\
kw\_private & binary & h1\_price\_mean & float & recent\_available & float \\
kw\_near & binary & h1\_price\_median & float & recent\_price\_mean & float \\
kw\_cozy & binary & h1\_price\_std & float & recent\_price\_std & float \\
kw\_view & binary & h1\_price\_nunique & integer & recent\_booking\_rate & float \\
kw\_modern & binary & h1\_price\_min & float & recent\_available\_rate & float \\
title\_len & integer & h1\_price\_max & float & price\_change\_q2\_q1 & float \\
title\_words & integer & h1\_price\_q25 & float & booking\_rate\_change\_q2\_q1 & float \\
geo\_cluster & str & h1\_price\_q75 & float &  &  \\
q1\_days\_total & integer & h1\_booking\_rate & float &  &  \\
\bottomrule
\end{tabular}
\end{table*}

\end{document}
